\documentclass[../math2.tex]{subfiles}
\begin{document}

\chapter{Puissances de nombres entiers}
\section{Pour se rappeler \dots}
%===============================================================================
%===============================================================================
\begin{exercise}
	Coche la ou les bonne(s) réponse(s).\par
	\begin{tasks}
		\task L'aire d'un carré de côté $2a$ est
		\begin{qcm}
			\choice $8a$
			\choice $2a^2$
			\choice $4a^2$
		\end{qcm}

		\task Le volume d'un cube d'arête $a$ est
		\begin{qcm}
			\choice $3a$
			\choice $a^2 \cdot a$
			\choice $a^3$
		\end{qcm}

		\task $20^0$ vaut
		\begin{qcm}
			\choice $1$
			\choice $20$
			\choice $0$
		\end{qcm}

		\task Le carré de $13$ est
		\begin{qcm}
			\choice $121$
			\choice $144$
			\choice $169$
		\end{qcm}

		\task Une puissance paire d'un nombre négatif est un nombre
		\begin{qcm}
			\choice négatif
			\choice positif
			\choice sans avis
		\end{qcm}
	\end{tasks}
\end{exercise}
%===============================================================================
\begin{solution}
	Coche la bonne réponse.\par
	\begin{tasks}
		\task L'aire d'un carré de côté $2a$ est
		\begin{qcm}
			\choice $8a$
			\choice $2a^2$
			\correct $4a^2$
		\end{qcm}

		\task Le volume d'un cube d'arête $a$ est
		\begin{qcm}
			\choice $3a$
			\choice $a^2 \cdot a$
			\correct $a^3$
		\end{qcm}

		\task Le carré de $13$ est
		\begin{qcm}
			\choice $121$
			\choice $144$
			\correct $169$
		\end{qcm}

		\task $\dfrac{2^3\cdot 3\cdot 5^2}{5\cdot 3^2}$ équivaut à
		\begin{qcm}
			\correct $2^3\cdot 3\cdot 5$
			\choice $\dfrac{2^3\cdot5}{3}$
			\choice $2^3$
		\end{qcm}

		\task Une puissance paire d'un nombre négatif est un nombre
		\begin{qcm}
			\choice négatif
			\correct positif
			\choice sans avis
		\end{qcm}
	\end{tasks}
\end{solution}
%===============================================================================
%===============================================================================
\begin{exercise}
	Complète les pointillés par le nombre qui convient.
	\begin{tasks}
		\task $6\cdot 6 \cdot 6 \cdot 6 = 6^{\dots}$
		\task $8=2^{\dots}$
		\task $(-5)\cdot(-5)\cdot(-5) = (-5)^{\dots}$
	\end{tasks}
\end{exercise}
%===============================================================================
\begin{solution}
	Remplace $n$ par le nombre qui convient.
	\begin{tasks}
		\task $6\cdot 6 \cdot 6 \cdot 6 = 6^4$
		\task $8=2^3$
		\task $(-5)\cdot(-5)\cdot(-5) = (-5)^3$
	\end{tasks}
\end{solution}
%===============================================================================
%===============================================================================
\newpage
\begin{exercise}
	Pour les puissances suivantes, fluote les bases,
	encadre les exposants et entoure la puissance. Ensuite, calcule.
	\begin{tasks}(2)
		\task  $8^2=$
		\task $(-3)^3=$
		\task $-2^4=$
		\task $5^3=$
	\end{tasks}
\end{exercise}
%===============================================================================
%===============================================================================
\begin{exercise}
	Calcule en écrivant les étapes de ton raisonnement.
	\begin{tasks}[after-item-skip=5em, after-skip=5em](2)
		\task $7^2 - 1^3=$
		\task $8\cdot(3-5)^2+4=$
		\task $3+2^2\cdot4^2=$
		\task $\left( 3^2 + 1 \right)^2 \cdot (-3)^2 + 5=$
	\end{tasks}
\end{exercise}

\newpage
\section{Produit de puissances de même base}
Voici une liste de produits de puissances. Décompose-les, si besoin, sous forme de produits et écris ensuite la puissance de $5$ correspondante.

\exemple
$\begin{aligned}[t]
		5^3 \cdot 5^7 & = \overbrace{5 \cdot 5 \cdot 5}^{5^3} \cdot \overbrace{5 \cdot 5 \cdot 5 \cdot 5 \cdot 5 \cdot 5 \cdot 5}^{5^7} \\
		              & = 5^{10}
	\end{aligned}$

%===============================================================================
%===============================================================================
\begin{tasks}[after-item-skip=2em, after-skip=2em](2)
	\task $5^3 \cdot 5=$
	\task $5^2 \cdot 5^4=$
	\task $5^2 \cdot 5^3=$
	\task $5^5 \cdot 5=$
	\task $5^3 \cdot 5^2 \cdot 5^2=$
	\task $5^{10} \cdot 5^2 \cdot 5=$
\end{tasks}
En observant la liste ci-dessus, écris $17^{10}\cdot17^{26}$ comme une puissance de $17$ : \dotfill

\adf[Pour réduire un produit de puissances de même base,]{2}
\begin{center}
	\mybox{$a^n\cdot a^p=\hspace*{1 in}$}
\end{center}

\begin{exercise}
	Réduis les produits suivants sous la forme $a^n$.
	\begin{tasks} (2)
		\task $7^4\cdot7=$
		\task $(-6)^3\cdot(-6)^5=$
		\task $3^9\cdot 3^{-5}=$
		\task $8^{-1}\cdot 8^3\cdot 8^2=$
		\task $12\cdot 12\cdot 12^3=$
		\task $\left(\dfrac{3}{2}\right)^7\cdot\left(\dfrac{3}{2}\right)^{-1}=$
	\end{tasks}
\end{exercise}

\begin{exercise}
	Complète les égalités suivantes.
	\begin{tasks} (2)
		\task $6\cdot6^{10}\cdot6^3\cdot \ldots \cdot 6=6^{20}$
		\task $11\cdot11^4\cdot 11^7=11^{\ldots}$
		\task $(-3)^3\cdot(-3)^4\cdot \ldots \cdot (-3)=(-3)^{13}$
		\task $\left(\dfrac{2}{5}\right)^6\cdot\left(\dfrac{2}{5}\right)^2 \cdot \dfrac{2}{5} \cdot \ldots = \left(\dfrac{2}{5}\right)^{20}$
	\end{tasks}
\end{exercise}

\section{Puissance d'un produit}
Tanguy et Johanna se retrouvent devant un énoncé qu'ils n'ont jamais vu : $(2\cdot3)^2$

Voici ce que chacun effectue comme calcul:

\parbox{.5\linewidth}{%
	\centering
	Méthode de Tanguy :
	\begin{align*}
		(2\cdot 3)^2
		 & = (6)^2 \\
		 & = 6^2   \\
		 & = 36
	\end{align*}}
\parbox{.5\linewidth}{%
	\centering
	Méthode de Johanna :
	\begin{align*}
		(2\cdot 3)^2
		 & = 2^2\cdot3^2 \\
		 & = 4\cdot 9    \\
		 & = 36
	\end{align*}
}

Ils sont étonnés de voir qu'ils n'ont pas utilisé le même procédé, mais que leur résultat est identique !

L'énoncé suivant est algébrique : $(m\cdot p)^b$ . Le professeur leur demande d'écrire le même calcul sans parenthèses.

Johanna s'écrie \og J'ai trouvé !\fg{} et écris : $(m\cdot p)^b=$\dotfill

\adf[Explique comment procéder pour élever un produit à une puissance :]{2}

\begin{center}
	\mybox{$(a\cdot b)^p=\hspace{1in}$}
\end{center}


\begin{exercise}
	\begin{tasks} (2)
		\task $(2\cdot11)^3=\dots^3\cdot\dots^3$
		\task $14^3=2^{\dots}\cdot7^{\dots}$
		\task $(2\cdot5)^8=10^{\dots}$
		\task $(e\cdot t)^7=$
		\task $(-2gh)^3=$
		\task $(3\cdot3\cdot7)^3=3^{\dots}\cdot 7^{\dots}$
		\task $2^{\dots} \cdot 5^{\dots}\cdot 49 = 10^4\cdot 7^{\dots}$
		\task $3^{\dots}\cdot5^{\dots}\cdot7^{\dots}=21^2\cdot125$
		\task $(5ab)^2=$
		\task $(3\cdot2\cdot11)^2=6^{\dots}\cdot121$
	\end{tasks}
\end{exercise}

\pagebreak
\section{Puissance d'une puissance}
Barre les fausses égalités.
\begin{tasks} [label=\,, after-item-skip=1em, after-skip=1em] (3)
	\task $(3^2)^3=3^5$
	\task $(3^2)^3=3^2\cdot 3^2\cdot 3^2$
	\task $ (3^2)^3=3^6$
	\task $(3^3)^2=3^5$
	\task $(3^3)^2=3^3\cdot3^3$
	\task $(3^3)^2=3^6$
\end{tasks}

\adf[Explique comment procéder pour élever une puissance à une puissance :]{2}

\begin{center}
	\mybox{$(a^m)^p=\hspace{1 in}$}
\end{center}

\begin{exercise}
	Réduis les produits suivants sous la forme $a^n$.
	\begin{tasks} (3)
		\task $(5^2)^4=$
		\task $(a^6)^3=$
		\task $\left((-2)^6\right)^2=$
		\task $\left(b^2\cdot b\right)^3=$
		\task $\left(2\cdot 2^7\right)^3=$
		\task $\left(13^4\cdot13^2\right)^2=$
	\end{tasks}
\end{exercise}

\begin{exercise}
	Complète les égalités suivantes.
	\begin{tasks} (2)
		\task $6^{14}=(6^2)^{\dots}$
		\task $2^7\cdot(2^3)^2=2^{\dots}$
		\task $2^{\dots}\cdot 3^{\dots}=(6^2)^2$
		\task $\left(2\cdot 3^5\right)^{\dots}=2^3\cdot 3^{\dots}$
		\task $\left(2a^3\right)^2=\dots a^{\dots}$
		\task $36a^6b^8=\left(6a^{\dots}b^{\dots}\right)^{\dots}$
	\end{tasks}
\end{exercise}
\pagebreak
\section{Quotient de puissances de même base}
Nous savons qu'une fraction représente le quotient du numérateur par le dénominateur.

Écris les quotients suivants sous forme soit d'une pusisance de 2, soit d'un produit d'une puissance de 2 et d'une puissance de 5.
\begin{tasks}[after-item-skip=4em, after-skip=4em] (2)
	\task $\dfrac{2^5}{2}=$
	\task $\dfrac{2^5}{2^2}=$
	\task $\dfrac{2^5}{2^3}=$
	\task $\dfrac{2^5}{2^4}=$
	\task $\dfrac{2^3\cdot 5^3}{2\cdot 5}=$
	\task $\dfrac{2^7\cdot 5^3}{2^4\cdot 5^2}=$
	\task $\dfrac{2^8\cdot 5^6}{2^8\cdot 5^6}=$
	\task $\dfrac{5^3\cdot 2^4}{2^3\cdot 5^2}=$
\end{tasks}

\adf[Comment procéder pour déterminer le quotient de puissances de même base?]{2}

\begin{center}
	\mybox{$\dfrac{a^p}{a^m}=\hspace{1in}$}
\end{center}

\begin{exercise}
	Écris sous la forme $a^n$.
	\begin{tasks} (2)
		\task $\dfrac{6^7}{6^5}=$
		\task $\dfrac{19^9}{19^3}=$
		\task $\dfrac{(-8)^7}{(-8)^3}=$
		%%	\task $\dfrac{(\frac{3}{4})^}10}{(\frac{3}{4})^8}$
		\task $\dfrac{\left(\dfrac{3}{4}\right)^{10}}{\left(\dfrac{3}{4}\right)^8}=$
	\end{tasks}
\end{exercise}

\section{Puissance d'un quotient}
\adf[Lorsqu'on élève une fraction à une puissance, on ]{1}

\begin{center}
	\mybox{$\left(\dfrac{a}{b}\right)^m=\hspace{1in}$}
\end{center}

\begin{exercise}
	Effectue. Veille à ce que ton résultat soir irréductible.
	\begin{tasks}(2)
		\task $\left(\dfrac{8}{3}\right)^2=$
		\task $\left(\dfrac{-3}{4}\right)^3=$
		\task $\left(\dfrac{4}{16}\right)^3=$
		\task $\left(\dfrac{1}{-2}\right)^4=$
	\end{tasks}
\end{exercise}

\section{Synthèse sur les propriétés des puissances}

 {\renewcommand{\arraystretch}{2}
  \begin{tabularx}{\linewidth}{|p{7cm}|X|X|}
	  \hline
	  \centering\textbf{Nom des propriétés} & \centering\textbf{Propriétés} & \centering\arraybackslash\textbf{Exemples} \\
	  \hline
	  Produit de puissances de même base    &                               &                                            \\
	  \hline
	  Puissance d'une puissance             &                               &                                            \\
	  \hline
	  Puissance d'un produit                &                               &                                            \\
	  \hline
	  Quotient de puissance de même base    &                               &                                            \\
	  \hline
	  Puissance d'un quotient               &                               &                                            \\
	  \hline
	  Puissance dont l'exposant vaut $0$    &                               &                                            \\
	  \hline
	  Puissance dont l'exposant vaut $1$    &                               &                                            \\
	  \hline
	  Puissance dont la base vaut $0$       &                               &                                            \\
	  \hline
	  Puissance dont la base vaut $1$       &                               &                                            \\
	  \hline
  \end{tabularx}}
\pagebreak
\begin{exercise}
	Complète les égalités suivantes.
	\begin{tasks} (2)
		\task $2^{\dots}\cdot 8^{\dots}=16^5$
		\task $\left(\dotfillbox [1cm] \right)^2=16x^2y^6$
		\task $\left(-6^4\right)^4=$
		\task $9^{\ldots} = 1$
		\task $3^3\cdot 3^9\cdot 3^{-2}\cdot 3^{\dots}=3^{15}$
		\task $\left(17^{\dots}\cdot 3^{\dots}\right)^3\cdot 9=3^{17}\cdot 17^6$
		\task $\left(9a^2b^5\right)^2=\ldots \cdot a^{\dots} \cdot b^{\dots}$
		\task $2^3\cdot 6^3=12^{\dots}$
		\task $\ldots^{1} = 12$
		\task $\left(5\cdot 2\cdot 11\right)^2=10^{\dots}\cdot 121$
		\task $3^2\cdot 9^2=3^{\dots}$
		\task $\left(\dfrac{6}{3}\right)^3=$
	\end{tasks}
\end{exercise}

\begin{exercise}
	Repère la seule expression qui correspond exactement à la proposition. Justifie ensuite en donnant le nom de la ou des propriété(s) correspondante(s).
	\begin{tasks}
		\task $(3x^3)^2$
		\begin{qcm}[4]
			\choice $9x^5$
			\choice $3x^6$
			\choice $9x^6$
			\choice $6x^5$
		\end{qcm}
		Propriété(s) : \dotfill
		\task $2x^2\cdot 3x^3$
		\begin{qcm}[4]
			\choice $36x^5$
			\choice $6x^5$
			\choice $108x^5$
			\choice $6x^6$
		\end{qcm}
		Propriété(s) : \dotfill
		\task $(2x)^2\cdot (2x)^3$
		\begin{qcm}[4]
			\choice $32x^6$
			\choice $25x^5$
			\choice $(2x)^5$
			\choice $4x^5$
		\end{qcm}
		Propriété(s) : \dotfill
		\task $3x^3\cdot 2x^3$
		\begin{qcm}[4]
			\choice $5x^3$
			\choice $5x^6$
			\choice $6x^6$
			\choice $6x^9$
		\end{qcm}
		Propriété(s) : \dotfill
	\end{tasks}
\end{exercise}



% \begin{exercise}
% 	Simplifie les fractions suivantes. Laisse ta démarche visible.
% 	\begin{tasks} (2)
% 		\task $\dfrac{20^4\cdot 3^4}{3^3\cdot 10^3}$
% 		\task $\dfrac{18^2}{12^3}$
% 		\task $\dfrac{36\cdot 10^2}{30^2}$
% 	\end{tasks}
% \end{exercise}
% L'exercice précédent n'a pas de sens pour moi, on n'aura pas encore fait les fractions et l'année passée avec les premières, on ne l'aura pas fait... On devrait peut-être complètement l'enlever 

\begin{exercise}
	Écris les exposants manquants.
	\begin{tasks}(2)
		\task $24^9$ est le produit de $24^7$ par $24^{\dots}$
		\task Le double de $2^6$ est $2^{\dots}$
	\end{tasks}
\end{exercise}

\newpage
\section{Puissances de 10}

Complète le tableau suivant.

	{\renewcommand{\arraystretch}{1.5}
		\begin{tabularx}{\linewidth}{|X|X|}
			\hline
			\centering \textbf{Puissances de $10$} & \centering\arraybackslash \textbf{Noms usuels} \\\hline
			$10^{-3} = $                           &                                                \\\hline
			$10^{-2} = $                           &                                                \\\hline
			$10^{-1} = $                           &                                                \\\hline
			$10^0 = $                              &                                                \\\hline
			$10^1 = $                              &                                                \\\hline
			$10^2 = $                              &                                                \\\hline
			$10^3 = $                              &                                                \\\hline
			$10^4 = $                              &                                                \\\hline
			$10^5 = $                              &                                                \\\hline
			$10^6 = $                              &                                                \\\hline
			$10^9 = $                              &                                                \\\hline
		\end{tabularx}}
\bigbreak
Tu viens d'utiliser les puissances de 10 (positives ou négatives).

Quelle est donc la signification des puissances de 10?\\
Pour quel que soit l'entier positif $n$,
\begin{align*}
	10^n=\underbrace{100\dots0}_{\text{n zéros}} &  &
	10^{-n}=\underbrace{0,0\dots01}_{\text{n zéros}}
\end{align*}

\begin{tcolorbox}
	Note : les propriétés des puissances s'appliquent aussi aux exposants négatifs !

	\exemple $\left(10^{-2}\right)^3 = 10^{-6}$
\end{tcolorbox}

\newpage
\begin{exercise}
	Effectue.
	\begin{tasks} (2)
		\task $8\cdot 10^9=$
		\task $12\cdot 10^{-2}=$
		\task $0,00023\cdot 10^8=$
		\task $157\cdot 10^{-5}=$
		\task $4,5478\cdot 10^8=$
		\task $82,68\cdot 10^{-7}=$
		\task $150000000\cdot 10^{-9}=$
		\task $0,857\cdot 10^4=$
		\task $45000\cdot 10^{-6}=$
		\task $3,1415 \cdot 10^0=$
	\end{tasks}
\end{exercise}

\section{Notation scientifique}
Jean et Muriel ont classé correctement en deux colonnes les expressions qui leur ont été données par leur professeur : les notations scientifiques, et celles qui ne le sont pas.

En observant le tableau, peux-tu trouver les éléments qui caractérisent la notation scientifique.

\begin{center}
	\begin{tabularx}{.8\linewidth}{YY} %
		\toprule
		\textbf{Notations scientifiques} & \textbf{\sout{Notations scientifiques}} \\
		\midrule
		$3,2\cdot 10^3$                  & $123\cdot 2^3$                          \\
		$2\cdot 10^{-2}$                 & $0,05 - 10^{-3}$                        \\
		$3,.25\cdot 10^9$                & $32,4\cdot 3^9$                         \\
		$6,456 \cdot 10^{-1}$            & $0,68\cdot 10^{-6}$                     \\
		$1,2\cdot 10^2$                  & $92 + 10^{-1}$                          \\
		$3,056\cdot 10^{-12}$            & $83\cdot 10^2$                          \\
	\end{tabularx}
\end{center}

\adf{4}

\exemples
\begin{itemize}[label=$\bullet$]
	\item l'écriture scientifique du nombre $456 789 000$ est \dotfill
	\item l'écriture scientifique du nombre $0,0000056$ est \dotfill
\end{itemize}

\begin{exercise}
	Parmi les notations suivantes, barre celles qui ne sont pas des notations scientifiques.
	\begin{tasks}[label=\,] (4)
		\task $0,5\cdot 10^9$
		\task $9,68\cdot 10^{-6}$
		\task $685\cdot 10^3$
		\task $1,23\cdot 20^3$
		\task $1,36\cdot 10^4$
		\task $25,6\cdot 10^{-3}$
		\task $123,123\cdot 10^2$
		\task $3,5^3$
	\end{tasks}
\end{exercise}

\begin{exercise}
	Transforme les nombres suivants en écriture décimale et en notation scientifique.

	\begin{tabularx}{\linewidth}{|Y|Y|Y|}
		\hline
		\textbf{Nombre}      & \textbf{Écriture décimale} & \textbf{Notation scientifique} \\\hline
		3 milliards          &                            &                                \\\hline
		$12,85\cdot 10^5$    &                            &                                \\\hline
		$0,036\cdot 10^7$    &                            &                                \\\hline
		$135,9\cdot 10^{-6}$ &                            &                                \\\hline
	\end{tabularx}
\end{exercise}

\begin{exercise}
	Classe les nombres suivants dans l'ordre croissant.
	\begin{tasks}[label=\,](3)
		\task $A=123000\cdot 10^{-4}$
		\task $B=123$
		\task $C=1,23\cdot 10^{-3}$
		\task $D=0,0000123\cdot 10^5$
		\task $E=1,23\cdot 10^3$
	\end{tasks}

	Classement :
\end{exercise}

\begin{exercise}
	Écris les produits suivants sous la forme d'une notation scientifique.
	\begin{tasks}[after-item-skip=8em, after-skip=8em]
		\task $\left(10^5\right)^2\cdot 5\cdot 10^{-2}\cdot 4\cdot \left(10^2\right)^4=$
		\task $7\cdot 10^{-9}\cdot 10^{-1}\cdot \left(10^7\right)^3\cdot 8=$
		\task $12\cdot 10^5\cdot 9\cdot 10^{-2}\cdot \left(10^8\right)^3=$
	\end{tasks}
\end{exercise}

\section{Exercices mélangés}
%Là peut-être que nos numéros d'exercices pourraient reprendre à 1? 

\begin{exercise}
	Vrai ou faux? Justifie dans tous les cas.
	\begin{tasks}
		\task $(a)^3$ sera toujours positif.
		\task $\left(\dfrac{3a}{bc}\right)^3 =\dfrac{3a^3}{b^3c^3}$
		\task $(-2)^2=-2^2$
		\task $-5a^3\cdot 6a= -30a^3$
		\task Vanessa prétend que le cube de $10^2$ est égal au carré de $10^3$.
		\task Nathan a écrit $(10^3)^2=10^9$.
		\task Nina pense que le cube d'une somme de deux nombres est égal à la somme des cubes de ces nombres.
		\task Karim pense que le cube d'un produit de deux nombres est égal au produit des cubes de ces nombres.
		\task Marie pense que le tiers de $3^{12}$ est $3^4$.
	\end{tasks}
\end{exercise}

\begin{exercise}
	Écris les produits suivants sous forme réduite et sans parenthèse.
	\begin{tasks} (2)
		\task $2x^2\cdot 3x$
		\task $5a^2\cdot a\cdot 3a^3$
		\task $\dfrac{(2ab)^2}{2a}$
		\task $(2x)^3\cdot (2x)^3$
		\task $\dfrac{\left(x^4\cdot x^6\right)}{x^2}$
		\task $(2x)^2\cdot xy^2$
		\task $\dfrac{a^7b^9\cdot a}{a^7\cdot b^3}$
		\task $10(2a^2\cdot 5a^8)$
		\task $5a^2\cdot 2a\cdot 3b\cdot 2b^2$
		\task $(2x^2)^3$
		\task $(3ab)^3$
		\task $(3xy)^2\cdot 2x^3y^6$
		\task $10\left(3x^2\cdot (x^3)^2\right)$
		\task $\dfrac{x^3\cdot x^6}{x^2\cdot x^2}$
	\end{tasks}
\end{exercise}

\begin{exercise}
	Quel est le nom de la propriété utilisée dans chacune des égalités suivantes?
	\begin{tasks} (2)
		\task $\left(\dfrac{5}{2}\right)^3=\dfrac{5^3}{2^3}$
		\task $\dfrac{8^7}{8^6}=8$
		\task $(5\cdot 9)^6=5^6\cdot 9^6$
		\task $x^6\cdot x=x^7$
	\end{tasks}
\end{exercise}

\begin{exercise}
	Cite la ou les propriété(s) utilisées dans les égalités suivantes.
	\begin{tasks}
		\task $(2\cdot 10)^3=8\cdot 1000$
		\task $x^8\cdot (x^2)^3=x^{14}$
	\end{tasks}
\end{exercise}

\begin{exercise}
	Écris le résultat sous forme de puissance d'un seul nombre.
	\begin{tasks} (3)
		\task $7^{25}\cdot 7^{31}$
		\task $(-9)^{16}\cdot (-9)^7$
		\task $\dfrac{5^{37}}{5^{25}}$
		\task $6^3\cdot (6^3)^2$
		\task $\dfrac{9^8}{9^6}$
		\task $(9^2)^6$
	\end{tasks}
\end{exercise}

\begin{exercise}
	Réduis au maximum les produits suivants.
	\begin{tasks}(2)
		\task $5a^4\cdot 2a\cdot 3b\cdot 2b^2$
		\task $3a^2\cdot a\cdot a^3b\cdot 3c^2$
		\task $x^2\cdot x\cdot x\cdot x^3\cdot x^7$
		\task $2^3\cdot (a^2)^3$
		\task $(3a^2b)^3$
		\task $(a^2)^4\cdot ab\cdot (b^3)^2$
		\task $2\cdot 3a^2\cdot a^3\cdot 2^2\cdot (a^3)^6$
		\task $(3xy)^2$
		\task $(3ab\cdot2ab)^3$
		\task $(5a^2b)^3\cdot (a^3b^2)^2$
		\task $(2ab^3)^2\cdot a\cdot b(a^2b)^3$
		\task $\dfrac{x^2\cdot x^5}{x\cdot x^2}$
	\end{tasks}
\end{exercise}

\begin{exercise}
	Trouve la valeur de $n$ dans les égalités suivantes et fluote les notations
	scientifiques.
	\begin{tasks} (2)
		\task $12500=1,25\cdot 10^n$
		\task $0,00125\cdot 10^n=125$
		\task $12500=n\cdot 10^3$
		\task $0,0001234=1,234\cdot 10^n$
		\task $104000000=n\cdot 10^6$
		\task $0,0001234=n\cdot 10^{-5}$
	\end{tasks}
\end{exercise}

\begin{exercise}
	Écris les nombres suivants en notation scientifique.
	\begin{tasks}(2)
		\task $250000000$
		\task $0,00005$
		\task $137\cdot 10^2$
		\task $35,69\cdot 10^{-5}$
		\task trente sept dixièmes
	\end{tasks}
\end{exercise}

\begin{exercise}
	Transforme ces mesures en mètres, et exprime le résultat en notation scientifique.
	\begin{tasks}(2)
		\task \SI{2500}{km}
		\task \SI{13}{cm}
		\task \SI{3}{mm}
		\task \SI{150000000}{km}
		\task \SI{250000}{cm}
		\task 364 milliards de km
	\end{tasks}
\end{exercise}

\begin{exercise}
	Classe les distances suivantes après les avoir écrites en notation scientifique.

	\begin{tabularx}{\linewidth}{YY}%
		\toprule
		\textbf{Planètes} & \textbf{Distances}    \\
		\midrule
		Terre - Soleil    & \SI{14960000}{km}     \\
		Terre - Pluton    & $5766\cdot 10^6$ km   \\
		Terre - Saturne   & 1,425 milliards de km \\
		\bottomrule
	\end{tabularx}
\end{exercise}

\begin{exercise}
	Calcule en utilisant les notations scientifiques. Écris ton résultat en notation scientifique.
	\begin{tasks}(3)
		\task $\np{60000}\cdot 0,002$
		\task $\np{0,000009}\cdot \np{300000000}$
		\task $\np{120000}\cdot \np{0,0002}$
	\end{tasks}
\end{exercise}

\begin{exercise}
	Un litre de yaourt contient 2 millions de bactéries. Combien de bactéries peut-on retrouver dans \SI{1250}{l} de yaourt? Écris ta réponses en notation scientifique.

\end{exercise}

\begin{exercise}
	Une population de microalgues se développe dans la marre du jardin. En un jour, elle se multiplie par 10.\\
	Si lundi, il y a 820 microalgues, combien se seront développées au bout de 7 jours ? \\
	Écris ton résultat en écriture scientifique.
\end{exercise}

\begin{exercise}
	La terre peut-être considérée comme une sphère de rayon de \SI{6400}{km}. L'aire $A$ d'une sphère de rayon $r$ est donnée par la formule
	\begin{equation*}
		A=4\pi r^2
	\end{equation*}
	Calcule, en \SI{}{km^2}, l'aire de la surface de la terre. Écris ton résultat en notation scientifique.


\end{exercise}

\begin{exercise}
	Le cerveau humain est composé de 100 milliards de neurones. A partir de 30 ans, ce nombre baisse d'environ 100 000 par jour. \\
	Donne les écritures décimale et scientifique du nombre de neurones qu'un humain a déjà perdu lorsqu'il atteint l'âge de 40 ans.

\end{exercise}
\end{document}